\section*{}
\section*{PREFACE TO THE FIRST EDITION}

Every aspect of this book was influenced by the desire to present calculus not
merely as a prelude to but as the first real encounter with mathematics. Since
the foundations of analysis provided the arena in which modern modes of math-
ematical thinking developed, calculus ought to be the place in which to expect,
rather than avoid, the strengthening of insight with logic. In addition to devel-
oping the students' intuition about the beautiful concepts of analysis, it is surely
equally important to persuade them that precision and rigor are neither deterrents
to intuition, nor ends in themselves, but the natural medium in which to formulate
and think about mathematical questions.
This goal implies a view of mathematics which, in a sense, the entire book
attempts to defend. No matter how well particular topics may be developed, the
goals of this book will be realized only if it succeeds as a whole. For this reason, it
would be of little value merely to list the topics covered, or to mention pedagogical
practices and other innovations. Even the cursory glance customarily bestowed on
new calculus texts will probably tell more than any such extended advertisement,
and teachers with strong feelings about particular aspects of calculus will know just
where to look to see if this book fulfills their requirements.
A few features do require explicit comment, however. Of the twenty-nine chap-
ters in the book, two (starred) chapters are optional, and the three chapters com-
prising Part V have been included only for the benefit of those students who might
want to examine on their own a construction of the real numbers. Moreover, the
appendices to Chapters 3 and 11 also contain optional material.
The order of the remaining chapters is intentionally quite inflexible, since the
purpose of the book is to present calculus as the evolution of one idea, not as a
collection of "topics." Since the most exciting concepts of calculus do not appear
until Part III, it should be pointed out that Parts I and II will probably require
less time than their length suggests—although the entire book covers a one-year
course, the chapters are not meant to be covered at any uniform rate. A rather
natural dividing point does occur between Parts II and III, so it is possible to
reach differentiation and integration even more quickly by treating Part II very
briefly, perhaps returning later for a more detailed treatment. This arrangement
corresponds to the traditional organization of most calculus courses, but I feel
that it will only diminish the value of the book for students who have seen a
small amount of calculus previously, and for bright students with a reasonable
background.
The problems have been designed with this particular audience in mind. They
range from straightforward, but not overly simple, exercises which develop basic
techniques and test understanding of concepts, to problems of considerable diffi-
culty and, I hope, of comparable interest. There are about 625 problems in all.
Those which emphasize manipulations usually contain many examples, numbered
vii